```latex
\documentclass[11pt]{article}
\usepackage[margin=1in]{geometry}
\usepackage{graphicx}
\usepackage{amsmath,amssymb}
\usepackage{booktabs}
\usepackage{hyperref}
\usepackage{authblk}
\usepackage{xcolor}
\usepackage{caption}

\title{\textbf{Autoscientist: A Multi-Agent Framework for Data Efficiency in Medical Imaging}}
\author{Autoscientist Collaboration}
\date{\today}

\begin{document}
\maketitle

\begin{abstract}
Data efficiency is a central concern in medical imaging, where labeled datasets are costly to assemble and prevalence varies dramatically across institutions. We present \textbf{Autoscientist}, a multi-agent framework that autonomously designs, executes, and interprets controlled experiments in medical machine learning. As a case study, Autoscientist evaluates whether \emph{prevalence matching}---downsampling training data so that the positive-class frequency matches that of the deployment distribution---improves out-of-distribution discrimination and calibration of pneumonia classifiers across four training-set sizes ($N \in \{1\text{k}, 5\text{k}, 25\text{k}, 100\text{k}\}$). Models are evaluated on two external chest X-ray datasets (NIH ChestX-ray14, PAD-CHEST). The headline interaction test does not reject the null hypothesis that matching moderates the effect of $N$ on AUROC ($F = 0.558$, permutation $p = 0.217$). Pairwise comparisons reveal that matching yields a statistically significant AUROC advantage only at $N=5{,}000$ ($\Delta = +0.066$, adj.\ $p = 0.021$), while at $N=25{,}000$ matching \emph{underperforms} the unmatched baseline ($\Delta = -0.031$) and produces 2.4$\times$ worse calibration. Near-chance AUROC values across conditions ($0.43$--$0.49$) indicate that domain shift, rather than prevalence composition, is the dominant bottleneck. We argue that autonomous multi-agent experimentation surfaces such non-monotonic, counterintuitive findings more efficiently than human-designed ablations alone.
\end{abstract}

\section{Introduction}

The promise of deep learning in medical imaging is repeatedly tempered by the brittleness of models when deployed across institutions with different patient populations, acquisition protocols, and disease prevalences \cite{zech2018variable, finlayson2021clinician}. Two intertwined challenges dominate: (i) \emph{data efficiency}---how to extract maximal generalization from limited labeled data---and (ii) \emph{distributional robustness}---how to ensure that performance on a development cohort transfers to a target deployment cohort.

A widely advocated mitigation is \emph{prevalence matching}: constructing the training set so that the marginal frequency of the positive class matches that of the deployment distribution \cite{larrazabal2020gender}. The intuition is that the implicit class prior learned by the model will align with the test distribution, improving both discrimination and calibration. However, prevalence matching is typically implemented by aggressive downsampling, which discards potentially informative examples. Whether the prior-alignment benefit outweighs the information loss is an empirical question whose answer likely depends on training-set size.

Answering such questions requires systematic experimentation across multiple sample sizes, matching strategies, model seeds, and evaluation cohorts. Manually designing and analyzing this combinatorial space is laborious and error-prone. We introduce \textbf{Autoscientist}, a multi-agent system in which specialized LLM-driven agents collaboratively (a) formulate hypotheses, (b) write training and evaluation code, (c) execute experiments on managed compute, and (d) statistically interpret results. As a representative case study, we use Autoscientist to investigate the prevalence-matching hypothesis for pneumonia classification with cross-domain evaluation. Our contributions are:

\begin{itemize}
    \item A multi-agent framework that autonomously executes a 2 (matching condition) $\times$ 4 (training size) factorial design with appropriate statistical analysis (linear mixed effects interaction test, Holm-corrected pairwise comparisons, permutation inference).
    \item An empirical finding that prevalence matching does \emph{not} consistently improve cross-domain AUROC, with a single isolated significant benefit at $N=5{,}000$ and a counterintuitive reversal at $N=25{,}000$.
    \item Evidence that calibration error is highly unstable for matched models at intermediate $N$, raising concerns about deploying naively prevalence-matched models.
\end{itemize}

\section{Methods}

\subsection{Autoscientist Framework}
Autoscientist orchestrates four cooperating agents: (1) a \emph{Hypothesis} agent that translates a high-level research question into a falsifiable statistical model; (2) an \emph{Engineering} agent that implements data loaders, training loops, and evaluation pipelines; (3) an \emph{Execution} agent that schedules and monitors training runs across multiple seeds; and (4) a \emph{Statistics} agent that performs preregistered analyses and produces the figures and tables. Inter-agent communication is mediated by structured handoff documents containing code artifacts, intermediate outputs, and audit trails.

\subsection{Case Study: Prevalence Matching for Pneumonia Classification}
\paragraph{Data.} Training data are drawn from a large publicly available chest X-ray corpus with binary pneumonia labels. Two external test sets are used: NIH ChestX-ray14 \cite{wang2017chestxray} and PAD-CHEST \cite{bustos2020padchest}. Each test set's empirical pneumonia prevalence is computed and used as the target for matching.

\paragraph{Conditions.} We vary training set size $N \in \{1{,}000, 5{,}000, 25{,}000, 100{,}000\}$ and matching condition $\in \{\text{matched}, \text{unmatched}\}$. In the \emph{matched} condition, negatives are subsampled so that the training-set positive prevalence equals the (average) test-set prevalence; in the \emph{unmatched} condition, the natural prevalence is preserved.

\paragraph{Model and training.} A standard ResNet-50 backbone is fine-tuned with cross-entropy loss for a fixed compute budget per condition. Each $(N, \text{condition})$ cell is replicated across multiple seeds.

\paragraph{Evaluation.} Primary metric: AUROC on the pooled NIH + PAD-CHEST test set. Secondary metrics: Expected Calibration Error (ECE) with 15 equal-width bins, and Brier score. The generalization gap is measured as the difference in AUROC between NIH and PAD-CHEST.

\paragraph{Statistical analysis.} The primary hypothesis is tested via a linear mixed-effects model with fixed effects for $\log N$, matching condition, and their interaction, with random intercepts per seed. The interaction term is the preregistered headline test; its significance is assessed by a permutation test ($10{,}000$ shuffles) over the condition labels. Pairwise comparisons at each $N$ use Holm correction across the four contrasts.

\section{Results}

\subsection{No Overall Effect of Matching}
The preregistered interaction test fails to reject the null hypothesis that matching moderates the effect of training-set size on AUROC: $F = 0.558$, permutation $p = 0.217$. There is therefore no overall evidence that prevalence matching offers a consistent benefit as a function of $N$.

\subsection{Pairwise AUROC}
Table~\ref{tab:auroc} summarizes pairwise contrasts. Matching produces a statistically significant gain only at $N=5{,}000$ ($\Delta = +0.066$, Holm-adj.\ $p = 0.021$). At $N=25{,}000$ the direction of the effect reverses: matched models underperform unmatched ones by 3.1 percentage points (adj.\ $p = 0.669$, not significant). At $N=1{,}000$ and $N=100{,}000$ matched models exhibit small, non-significant advantages.

\begin{table}[h]
\centering
\caption{Mean AUROC by training-set size and matching condition, with Holm-adjusted $p$-values from pairwise contrasts.}
\label{tab:auroc}
\begin{tabular}{rcccc}
\toprule
$N$ & Matched & Unmatched & $\Delta$ (M$-$U) & Adj.\ $p$ \\
\midrule
1{,}000   & 0.486 & 0.427 & $+0.060$ & 0.669 \\
5{,}000   & 0.491 & 0.425 & $+0.066$ & \textbf{0.021} \\
25{,}000  & 0.446 & 0.477 & $-0.031$ & 0.669 \\
100{,}000 & 0.465 & 0.432 & $+0.033$ & 0.669 \\
\bottomrule
\end{tabular}
\end{table}

All AUROC values cluster near $0.45$--$0.49$, near the $0.5$ chance baseline, indicating substantial domain shift that neither condition fully overcomes. The measured generalization gap between NIH and PAD-CHEST in a representative run is $0.068$ AUROC, comparable in magnitude to any matching-induced effect.

\subsection{Calibration}
Calibration results (Table~\ref{tab:cal}) reveal a more dramatic and concerning pattern. At $N=1{,}000$ and $N=25{,}000$, matched models are markedly \emph{worse} calibrated than unmatched models (ECE $\approx 0.31$ vs.\ $0.10$ at $N=1$k; $0.33$ vs.\ $0.13$ at $N=25$k). At $N=5{,}000$ and $N=100{,}000$, matched models are slightly better calibrated. None of the calibration differences survive multiple-testing correction, but the effect sizes at the two flagged sample sizes are large. The non-monotonic ``bad--good--bad--good'' pattern for matched models suggests training instability induced by aggressive downsampling.

\begin{table}[h]
\centering
\caption{Expected Calibration Error (ECE) and Brier score by condition.}
\label{tab:cal}
\begin{tabular}{rcccc}
\toprule
$N$ & ECE (M) & ECE (U) & Brier (M) & Brier (U) \\
\midrule
1{,}000   & 0.313 & 0.097 & 0.210 & 0.067 \\
5{,}000   & 0.088 & 0.118 & 0.060 & 0.073 \\
25{,}000  & 0.327 & 0.134 & 0.248 & 0.111 \\
100{,}000 & 0.048 & 0.067 & 0.048 & 0.063 \\
\bottomrule
\end{tabular}
\end{table}

\section{Discussion}

\paragraph{Prevalence matching is not a free lunch.} Our preregistered interaction test does not support the hypothesis that prevalence matching consistently improves cross-domain discrimination. The only positive significant result occurs at $N=5{,}000$, an intermediate regime where the matched training set is plausibly large enough to be informative yet small enough that prior alignment confers a marginal benefit. The reversal at $N=25{,}000$ is the most striking observation: matched models lose both AUROC and---more sharply---calibration relative to their unmatched counterparts. A plausible mechanism is that constructing a 25k matched set from a much larger pool requires discarding a large fraction of informative negatives, distorting the effective negative distribution while simultaneously presenting the model with a near-balanced prior that is misaligned with the eventual test mixture.